\documentclass[sigconf,nonacm]{acmart}
\usepackage{booktabs}
\usepackage{multirow}
\usepackage{xcolor}
\usepackage{xurl}
\settopmatter{printacmref=false}
\renewcommand\footnotetextcopyrightpermission[1]{}
\pagestyle{plain}

\newcommand{\code}[1]{\texttt{#1}}

\setlength{\emergencystretch}{2em}
\begin{document}

\title{Traversing a Learning-Progression Knowledge Graph to Diagnose Novice Code-Comprehension Reasoning}
\subtitle{Mid-Term Progress Report}

\author{Pradhyumn Nair}
\affiliation{%
  \institution{Department of Computer Science, Georgia Southern University}
  \city{Statesboro}
  \state{GA}
  \country{USA}}
\email{pn02774@georgiasouthern.edu}

\begin{abstract}
A correct answer to a code-tracing question is weak evidence of correct reasoning, and a wrong answer
rarely says which piece of understanding is missing. We study whether a knowledge graph that encodes
a learning progression can make that diagnosis explicit and traceable. The graph covers Chapter~1 of an
introductory Java textbook: 173 concepts joined by 241 prerequisite and parallel edges that each carry a
written \emph{edge logic}, 296 literature-sourced misconceptions and errors attached through 569 edges
that pair the wrong rule with its correct twin, and five level nodes (L0--L4) per concept whose 3{,}079
``ways'' state what a student at that level says. Our pipeline has two stages. Stage~1, implemented and
tested, reads only the program text, matches recognition patterns stored on the concept nodes, and
extracts the subgraph of concepts, edges, misconceptions and levels relevant to the item. Stage~2, now
in design, adapts Think-on-Graph: a language model traverses that subgraph with beam search, pruning
edges against the student's clauses, and learning-progression rules turn the matched edges into a level
per concept and a root cause. On 10 edge-case scenarios and one reference case, Stage~1 extracted a subgraph
for every item: 2 to 26 concepts (9.6 on average) with their relation edges, misconceptions and L0--L4 levels,
cutting the misconception edges a reasoner must consider from 569 to 63 on average. All 208 tests of the
extractor pass.
\end{abstract}

\keywords{knowledge graph, learning progression, code comprehension, misconceptions, Think-on-Graph,
large language models, computing education}

\maketitle

\section{Introduction}
Reading a program is a different skill from writing one, and it develops in stages
\cite{lister2004multinational,lister2006forest}. Novices hold specific, repeatable wrong ideas about what
the machine does: that a loop condition is watched continuously, that a variable holds an expression rather
than a value, that a local variable starts at zero \cite{pea1986bugs,duboulay1986difficulties,sorva2013notional}.
These gaps matter because a wrong mental model can still produce a right answer. On an open online course,
about one third of students could not explain their own working code \cite{lehtinen2021explain}, and in a
first-year Java course only 17\% of students held a sound model of reference assignment \cite{ma2007viability}.

Grading the output of a tracing question therefore measures the wrong thing. An instructor, or an automated
tutor choosing the next exercise, needs to know \emph{which} piece of understanding is missing and \emph{how
far along} the learner is. Computing education has scales of depth (SOLO \cite{biggs1982solo}, the Block
Model \cite{schulte2008block}) and large catalogues of misconceptions \cite{qian2017misconceptions,chiodini2021progmiscon},
and science education has \emph{learning progressions}: ordered, validated descriptions of how reasoning about
one idea develops \cite{krajcik2023progressions,wilson2009measuring}. Our scoping review of 240 studies
\cite{nair2026scoping} found that these three ingredients (code-comprehension assessment, validated ordered
scales, and language-model scoring) have each been paired, but not yet joined.

Large language models (LLMs) can now read free-text explanations and grade them at scale
\cite{fowler2021autograding,zhang2025solo,prather2023robots}, but their verdicts are hard to trace to a rule, and LLM judges show known biases \cite{zheng2023judging}.
Recent work on reasoning over knowledge graphs offers a remedy: Think-on-Graph lets an LLM explore a graph
step by step and answer from the paths it found \cite{sun2024tog}, and Reasoning on Graphs restricts the model to
paths that really exist in the graph \cite{luo2024rog}.

This project joins the two lines. We built a knowledge graph that encodes a learning progression for
Chapter~1 of Sedgewick and Wayne's \emph{Introduction to Programming in Java} \cite{sedgewick2017java}, with the
reasoning written on every edge, and we are building a pipeline that traverses it to diagnose a student's
explanation. The contributions reported at mid-term are:
\begin{itemize}
\item a learning-progression knowledge graph (KG) whose edges carry their logic, whose misconceptions carry a
correct twin, and whose level nodes carry the ways a student at each level reasons (Section~\ref{sec:data});
\item a tested, data-driven extractor that reads only program text and returns the relevant subgraph, with every
recognition pattern stored on the concept nodes rather than in code (Section~\ref{sec:stage1});
\item a design for adapting Think-on-Graph traversal and learning-progression rules to diagnosis, and preliminary
evidence on 10 scenarios (Sections~\ref{sec:stage2}--\ref{sec:results}).
\end{itemize}

\section{Research Questions}
\label{sec:rq}
\begin{description}
\item[RQ1] \emph{Subgraph recovery.} To what extent can the concepts, edges, misconceptions and level nodes
relevant to a code-comprehension item be recovered from the program text alone, using recognition patterns
stored on the KG's concept nodes? \emph{Measures:} size and connectedness of each subgraph, reduction of the
search space, and the extractor's test results.
\item[RQ2] \emph{Traversal.} How accurately does a Think-on-Graph traversal of the extracted subgraph align
each clause of a student's explanation with the edge or misconception it uses? \emph{Measures:} following
Hwang et al.\ \cite{hwang2025principal}, each edge and misconception of the subgraph is labelled \emph{correct},
\emph{incorrect} or \emph{not applied}, and we report knowledge precision, knowledge recall and F1 against expert
labels; exact-match accuracy of the held misconception and the root cause, as Hits@1 in ToG \cite{sun2024tog}.
\item[RQ3] \emph{Learning-progression logic.} How well do the level rules stored in the graph place a student
on L0--L4 for each concept, compared with expert labels? \emph{Measures:} exact-match accuracy of the level,
agreement with the acceptable set of levels, and weighted Cohen's $\kappa$.
\item[RQ4] \emph{Value of the graph.} Does grounding an LLM in the traversal improve diagnostic accuracy and
run-to-run consistency over the same LLM without the graph? \emph{Measures:} differences in the RQ2--RQ3 measures,
agreement across five reruns, and LLM calls per answer.
\end{description}
RQ1 is answered with preliminary results in this report; RQ2--RQ4 define the remaining work.

\section{Literature Review}
\label{sec:lit}
We organise prior work around the three ingredients our system combines, and the method that joins them.

\subsection{What novices get wrong, and how it is catalogued}
Early work identified language-independent ``bugs'' in novice thinking \cite{pea1986bugs} and the separate
difficulties novices face at once: the notional machine, notation, standard structures and pragmatics
\cite{duboulay1986difficulties}. Sorva's notional machines explain many of them \cite{sorva2013notional}, and Qian and
Lehman's review \cite{qian2017misconceptions} organises them into syntactic, conceptual and strategic
knowledge. Misconceptions have since been identified through formal interviews with students, which revealed four themes
\cite{kaczmarczyk2010identifying},
validated for Java toward a concept inventory \cite{caceffo2019java}, and curated into ProgMiscon, an online
inventory of 198 misconceptions \cite{chiodini2021progmiscon}. Learning-curve analysis separates the errors
practice fixes from those that persist \cite{rivers2016learning}. \emph{Gap:} these catalogues are flat lists.
They record that a misconception exists, but not which prerequisite relation it breaks or which level of
understanding it implies, which is what a diagnosis needs.

\subsection{Scales of depth for code comprehension}
The SOLO taxonomy grades an answer by how far its parts are tied together \cite{biggs1982solo}, and was
applied to novice code-reading responses \cite{lister2006forest}, building on a multinational tracing study \cite{lister2004multinational}.
The Block Model grades the reading itself, by the size of the program unit and the depth from text to purpose
\cite{schulte2008block}. Both order responses, not learners, and both depend on human raters whose agreement
must be checked. Knowledge tracing models a learner's mastery over time, from Bayesian knowledge tracing
\cite{corbett1994knowledge} to deep knowledge tracing \cite{piech2015deep}, but models correctness on items,
not the reasoning behind an answer. \emph{Gap:} no scale for code comprehension has been built from, and tested
against, the specific wrong ideas students hold.

\subsection{Learning progressions}
Science education describes how reasoning about one idea develops as a learning progression: ordered levels
drawn from real student answers \cite{krajcik2023progressions}. Wilson's construct-modeling approach ties each
level to assessment items \cite{wilson2009measuring}; Alonzo and Steedle built and tested a force-and-motion
progression with ordered multiple-choice items whose options map to levels \cite{alonzo2009force}; and Jin
et al.\ set out the evidence a progression's validity requires \cite{jin2019validation}. Automated scoring against
a progression is now feasible: machine scores can be compared directly with researchers' interpretations
\cite{jin2025automated}. \emph{Gap:} computing education has no validated progression for code comprehension,
and existing progressions are documents for human raters, not machine-readable structures that a system can
traverse.

\subsection{Language models for grading and tutoring}
Since 2022, LLMs have entered computing education as generators of material, explainers and tutors
\cite{prather2023robots}. CodeHelp showed guarded LLM help at course scale \cite{liffiton2023codehelp}, automatic
grading of ``explain in plain English'' answers reached accuracy close to trained teaching assistants
\cite{fowler2021autograding}, and LLMs now classify student programming answers onto SOLO levels
\cite{zhang2025solo}. Using LLMs as judges has known weaknesses: position, verbosity and self-enhancement biases, and weak grading of math and reasoning \cite{zheng2023judging}. \emph{Gap:} these systems apply a
general-purpose scale through a prompt. None constrains the model to a validated, ordered structure of the
domain, and none returns the specific prerequisite that failed.

\subsection{Reasoning over knowledge graphs with language models}
Graph-based methods let an LLM reason over explicit structure rather than from memory. Think-on-Graph (ToG)
treats the LLM as an agent that runs beam search on a knowledge graph. At each depth it lists the relations at
the current entities, prunes them, follows the kept ones, and asks itself whether the paths found so far
suffice \cite{sun2024tog}. It is training-free and set new best results on six of nine benchmarks with GPT-4.
Reasoning on Graphs (RoG) instead plans relation paths, retrieves only the paths that exist in the graph, and
reasons over them, which makes answers faithful and interpretable \cite{luo2024rog}. PathRAG shows that giving
the LLM key relational paths as ordered text, rather than a flat list of nodes and edges, improves answer
quality \cite{chen2026pathrag}. Principal-knowledge grounding checks reasoning against a set of required
knowledge units and labels each correct, incorrect or not applied \cite{hwang2025principal}, but its units are a
flat list without relations. \emph{Gap:} these methods answer factual questions over encyclopedic graphs. None
has been applied to diagnosing a student's reasoning, where the graph encodes prerequisites, misconceptions
and levels, and where the output is a level and a root cause, not an entity.

\subsection{Positioning}
Our work sits at the intersection the review identified \cite{nair2026scoping}. We make the learning progression
machine-readable as a knowledge graph, and use graph traversal, in the spirit of ToG and RoG, to tie an LLM's
reading of a student explanation to explicit edges, misconceptions and levels.

\section{Methodology}
\label{sec:method}
Figure~\ref{fig:pipeline} shows the pipeline. Stage~1 is implemented; Stage~2 is being built.

\begin{figure*}[t]
  \centering
  \includegraphics[width=0.95\textwidth]{figures/fig_pipeline.pdf}
  \caption{Pipeline. Stage~1 (implemented) reads only the code and extracts the relevant subgraph using
  patterns stored on the concept nodes. Stage~2 (planned) adapts Think-on-Graph to traverse that subgraph against
  the student's answer and applies the learning-progression rules.}
  \label{fig:pipeline}
\end{figure*}

\subsection{Dataset}
\label{sec:data}
\paragraph{Knowledge graph.} The KG (v3.9.5-reader4, JSON; also exported as a spreadsheet) covers Chapter~1 of
\cite{sedgewick2017java}: 6 sections, 36 topics and 173 concepts. Table~\ref{tab:kg} summarises it, and
Figure~\ref{fig:schema} shows its schema. Each prerequisite or parallel edge carries an \emph{edge logic}, a sentence
stating why the relation holds; for example, \code{loop\_control}$\to$\code{for\_loop}: ``The for header is exactly the
three control slots, and the body runs after the test and before the update on each pass.'' Each misconception
edge carries the wrong rule (\code{logic}) and the correct rule (\code{twin\_logic}). 294 of the 296 misconception and
error nodes cite literature, and their provenance is published work (167), the textbook (92), derived (28), a rule
dictionary (7) or the scenarios (2). Every concept has five level nodes, L0--L4, aligned with SOLO-style depth: L0 holds
a rule that contradicts Java; L1 names the construct or states its rule; L2 makes a correct prediction or trace;
L3 relates the concept to a neighbour; L4 diagnoses an error or predicts the effect of a change. Each level node is
reached by a \code{has\_level} edge listing its \emph{ways}, 3{,}079 in total; 736 of them link a specific misconception,
so a held misconception maps directly to a level. Each concept also has a \emph{reader} block with recognition
patterns; 166 concepts are visible in code (359 regular expressions), 7 are not (for example \code{testing} and
\code{flowchart}), and 21 are marked generic (for example \code{statement}).

\begin{table}[t]
\caption{The learning-progression knowledge graph (v3.9.5-reader4).}
\label{tab:kg}
\small
\begin{tabular}{@{}p{0.68\columnwidth}r@{}}
\toprule
Element & Count \\
\midrule
Sections / topics / concepts & 6 / 36 / 173 \\
Prerequisite / parallel edges (with edge logic) & 207 / 34 \\
Misconceptions / errors & 255 / 41 \\
Misconception and error edges (wrong + twin rule) & 569 \\
Level nodes (L0--L4 per concept) & 865 \\
Level ways / ways linked to a misconception & 3{,}079 / 736 \\
Concepts with code patterns / regexes & 166 / 359 \\
\bottomrule
\end{tabular}
\end{table}

\begin{figure}[t]
  \centering
  \includegraphics[width=\columnwidth]{figures/fig_schema.pdf}
  \caption{KG schema. A held misconception reaches a level through the way that names it.}
  \label{fig:schema}
\end{figure}

\paragraph{Scenarios.} We use the scenario file in the repository's \code{data/raw} folder: 10
edge-case scenarios (S001, S007, S008, S021, S026, S047, S052, S057, S134, S137) and one reference case (S000), drawn
from a Chapter~1 scenario set; each has code, a question and a student answer. They were chosen as hard cases:
self-corrections, correct answers by coincidence, non-native English. Stage~1 reads only the code; the question and
answer are inputs to Stage~2.


\paragraph{Preprocessing.} String and character literals and comments are blanked to spaces, keeping the
length so positions and line numbers stay valid. The code is split into statements on semicolons and braces only at parenthesis depth zero, so a \code{for} header stays whole. Shell-command fragments
(\code{java X < file}) are extracted for the I/O concepts. The student answer is split into clauses at link words
stored in the KG (\emph{because}, \emph{so}, \emph{but}, \emph{oh wait}, \ldots); this step belongs to Stage~2.

\subsection{Stage 1: subgraph extraction (implemented)}
\label{sec:stage1}
The extractor is data-driven. It loops over the concept nodes and runs each node's own patterns against the
code: regular expressions over the whole program, structural signals over the program or over each statement,
operator symbols, Java type names, and command patterns. A concept is tagged only when one of its own patterns
matches, and the code contains no concept identifiers. Generic concepts are kept only when a prerequisite or
parallel edge ties them to a specific tagged concept, or a decisive signal of theirs fired. The subgraph is then
read from the KG: every relation edge with both ends among the concepts, every misconception and error edge
from them, every \code{has\_level} edge with its ways, and the level--misconception links inside the ways.
Figure~\ref{fig:s007} shows the result for scenario S007.

\begin{figure*}[t]
  \centering
  \includegraphics[width=0.92\textwidth]{figures/fig_s007_crop.png}\\[2pt]
  \includegraphics[width=0.5\textwidth]{figures/fig_s007_legend.png}
  \caption{Stage~1 subgraph for S007 (nested loops; the student says ``the loop stops at 1''), produced by the
  repository's results script: 17 concept nodes, 14 relation edges, 141 misconception and error edges, 85 level nodes,
  in three disconnected pieces. Dark nodes were found by their own code patterns; light nodes are generic concepts kept
  because an edge ties them to a specific one.}
  \label{fig:s007}
\end{figure*}

\paragraph{Stage 1 on S007, step by step.} The S007 code is a \code{for} loop over \code{i} from 3 down to 0, with an
inner \code{for} loop that prints \code{j} with \code{printf} and a \code{println} after each outer pass.
Table~\ref{tab:s007tags} lists what tagged each specific concept. Every tag comes from a pattern stored on that concept's
own node. Two are structural signals rather than plain patterns: \code{loop\_control} is tagged by its \code{for\_header}
signal, which needs the whole header \code{(int i = 3; i >= 0; i--)}, and \code{nested\_loops} by its \code{nested\_for}
signal, which finds a second \code{for} inside the first loop's braces. Five generic concepts are kept because an edge
ties each to a specific one (for example \code{string\_literal}$\to$\code{println}). The KG then supplies 14 relation
edges among the 17 concepts, including \code{loop\_control}$\to$\code{for\_loop} and \code{for\_loop}$\to$\code{nested\_loops},
and the concepts fall into three pieces: the loop concepts, the output concepts, and \code{loop\_body} on its own.

\begin{table}[t]
\caption{What tagged each specific concept in S007 (line of the code where the match starts).}
\label{tab:s007tags}
\footnotesize
\begin{tabular}{@{}lll@{}}
\toprule
Concept & Tagged by & Line \\
\midrule
\code{for\_loop} & pattern \code{for (} & 1 \\
\code{loop\_control} & signal \code{for\_header} & 1 \\
\code{nested\_loops} & signal \code{nested\_for} & 1 \\
\code{loop\_variable}, \code{loop\_scope} & patterns on the \code{for} header & 1 \\
\code{loop\_body}, \code{control\_flow} & pattern \code{for (} / \code{if (} \ldots & 1 \\
\code{boolean\_expression} & pattern for a \code{for} test & 1 \\
\code{comparison\_operators} & pattern \code{>=}, \code{>} & 1 \\
\code{increment\_operators} & pattern \code{--} & 1 \\
\code{formatted\_output} & pattern \code{.printf(} & 3 \\
\code{println} & pattern \code{System.out.println(} & 5 \\
\bottomrule
\end{tabular}
\end{table}

\subsection{Stage 2: Think-on-Graph traversal and learning-progression logic (design)}
\label{sec:stage2}
We adapt ToG \cite{sun2024tog} from question answering to diagnosis, keeping its three phases.
\emph{Initialization:} the main concepts of the item are the topic entities. \emph{Exploration:} at each depth,
relation search lists the edges at the frontier within the subgraph (relation edges with their logic, misconception
edges with wrong and twin rules, and level ways), and relation pruning asks the LLM which of them the current
clause uses, keeping the top $N$; entity exploration follows the kept edges. \emph{Reasoning:} the LLM is asked
whether the paths found are sufficient to judge the clause; if not, exploration continues to depth $D_{\max}$.
Following ToG, $N=D_{\max}=3$, which bounds the cost at $2ND+D+1=22$ LLM calls per clause. Paths are given to
the LLM as ordered text, node, edge logic, node, as PathRAG recommends \cite{chen2026pathrag}, and the LLM sees
only edges that exist, as in RoG \cite{luo2024rog}. Figure~\ref{fig:tog} sketches the intended traversal for S007.

The learning-progression logic then turns matched edges into levels, using only rules stored in the KG.
(i)~A clause closer to a misconception's wrong rule than to its twin, and confirmed by a fact from the code (for S007,
the last pass is $i=0$), holds that misconception; the level is the L0 way that links it. (ii)~A relation edge whose
logic the clause states correctly is confirmed and supports L3 for its ends; an edge whose end holds a misconception
is broken. (iii)~Otherwise the level is the highest whose ways the clause matches, using the KG's evidence classes
(named and rule: L1, run: L2, relation: L3, diagnosis and hypothetical: L4). (iv)~The root cause is found by walking
up prerequisite edges from the main concept to the first broken edge.

\begin{figure}[!htb]
  \centering
  \includegraphics[width=0.62\columnwidth]{figures/fig_tog_s007.pdf}
  \caption{Intended Think-on-Graph traversal for clause c2 of S007 (design illustration, not a result).}
  \label{fig:tog}
\end{figure}

\paragraph{Worked example (design).} In S007 the student writes ``3 2 1 / 2 1 / 1 (no empty line at the end),
because the loop stops at 1.'' The clause splitter yields c1 (the output claim) and c2 (the reason, after
\emph{because}). Starting from \code{nested\_loops}, the traversal for c1 keeps \code{for\_loop}$\to$\code{nested\_loops},
whose logic (``the inner for runs its complete set of passes for every single pass of the outer for'') matches the
three rows; that edge is confirmed. For c2 it continues to \code{loop\_control}$\to$\code{for\_loop} and then to the
misconception edge \code{loop\_control}$\to$OffByOneBound, where the clause is closer to the wrong rule than to the twin.
The fact read from the loop header (test \code{i >= 0}, so the last pass is $i=0$) confirms that the student is wrong. The
rules then give \code{loop\_control} L0 through its linked way, mark \code{loop\_control}$\to$\code{for\_loop} broken, give
\code{nested\_loops} L0 because its output claim is contradicted, and name \code{loop\_control} as the root cause. This is the
diagnosis the implementation will be tested against.

\subsection{Training procedure}
Stage~1 needs no training: the patterns are written into the KG and tested. Stage~2 follows ToG in being
training-free: the LLM is used through prompts with a few worked examples, temperature~0 for the reasoning step,
and no fine-tuning. If prompting proves unreliable on near-identical wrong and twin rules, we will fine-tune a small
open model on pairs generated from the KG itself (wrong rule, twin rule, edge logic and level ways as labelled
statements), as RoG builds its training data from graph paths \cite{luo2024rog}; the 10 scenarios will be held
out for testing only. Hyperparameters are fixed in advance: $N=3$, $D_{\max}=3$, at most 800 characters between a
declaration and its use in bounded patterns. Stage~1 runs on a laptop CPU; Stage~2 calls a hosted LLM through its API.

\subsection{Evaluation}
Stage~1 is evaluated by its test suite and by the subgraphs it produces (RQ1). The test suite checks 372 snippets,
showing that each of the 166 code-visible concepts is tagged when it should be and not when it should not, plus core-rule
tests that patterns come only from the nodes. Stage~2 will be evaluated in two conditions: on the Stage~1 subgraph,
and on a subgraph built from expert-chosen concepts, which separates traversal errors from extraction errors. The
baselines are (B1) the same LLM given the code, answer and level definitions but no graph; (B2) a RoG-style single call
given the subgraph paths without iterative traversal; and (B3) a lexical TF-IDF judge in place of the LLM. Ablations
remove the prerequisite edges, the twin rules or the level ways. Consistency is measured over five reruns.

\paragraph{Output of a diagnosis.} For each student answer the system returns the overall level, a level per concept,
the misconceptions held, the root-cause concept, and for each verdict the evidence path in the KG (for S007:
\code{loop\_control}$\to$OffByOneBound$\to$\code{loop\_control@L0}). An answer is correct under exact match only if the
predicted level, misconception or root cause equals the expert label.

\paragraph{Result metric.} Hwang et al.\ assess a model's reasoning against the knowledge units a question requires,
labelling each unit \emph{correct}, \emph{incorrect} or \emph{not applied}, and report knowledge precision (of the units
used, the share used correctly), knowledge recall (of the units required, the share used correctly) and their F1
\cite{hwang2025principal}. We adopt this metric with the edges and misconceptions of the subgraph as the knowledge units:
a confirmed edge is \emph{correct}, a broken edge or a held misconception is \emph{incorrect}, and an edge the student
never uses is \emph{not applied}. For S007, \code{for\_loop}$\to$\code{nested\_loops} is correct and
\code{loop\_control}$\to$\code{for\_loop} is incorrect through OffByOneBound. Precision, recall and F1 are computed
between the system's labels and the expert's labels for each scenario.

\paragraph{Validating an LLM judge.} Expert labels are the reference. Where an LLM labels units at scale, we follow
Zheng et al., who validate an LLM judge by its agreement with human judgments and document its position, verbosity and
self-enhancement biases \cite{zheng2023judging}. We report the LLM judge's agreement with the expert labels (percent
agreement and Cohen's $\kappa$) on the 10 scenarios before using its labels, shuffle the order in which edges are
presented, use temperature~0, and use a different model family for judging than for reasoning.

\section{Results and Findings So Far}
\label{sec:results}
\paragraph{Stage~1 extracts a subgraph for every scenario.} Table~\ref{tab:rq1} and Table~\ref{tab:perscenario} report
RQ1. For all 10 scenarios and the reference case, the extractor found concept nodes from the code and read their
relation edges, misconceptions and level nodes out of the KG. Subgraphs range from 2 to 26 concepts (9.6 on average,
Figure~\ref{fig:coverage}). The repository's report prints ``fully covered: 10 / 10'', and all 208 tests pass in 2.3
seconds.

\begin{table}[t]
\caption{RQ1: Stage~1 subgraphs over the 10 scenarios (S000 reference case excluded).}
\label{tab:rq1}
\small
\begin{tabular}{@{}p{0.50\columnwidth}ccc@{}}
\toprule
Per subgraph & Mean & Min & Max \\
\midrule
Concept nodes & 9.6 & 2 & 26 \\
Generic concepts kept & 2.5 & 1 & 5 \\
Relation edges & 8.2 & 1 & 22 \\
Misconception and error edges & 63.0 & 7 & 167 \\
Level nodes & 48.0 & 10 & 130 \\
Level--misconception links & 69.0 & 11 & 193 \\
Disconnected pieces & 2.0 & 1 & 4 \\
\bottomrule
\end{tabular}
\end{table}

\begin{figure}[t]
  \centering
  \includegraphics[width=\columnwidth]{figures/fig_coverage.pdf}
  \caption{Concept nodes in each Stage~1 subgraph: found by their own pattern (dark) and generic concepts kept because an
  edge ties them to a specific one (light).}
  \label{fig:coverage}
\end{figure}

\paragraph{The graph narrows the search.} A subgraph holds 63 of the KG's 569 misconception and error edges on average
(Figure~\ref{fig:searchspace}), about 11\% of the graph, so a reasoner works with a small, relevant set rather than the
whole catalogue. Scenarios with short code (S026, S047, S134) give subgraphs of two concepts; longer code with loops,
arrays and output (S007, S008, S057) gives 17 to 26 concepts.

\begin{figure}[t]
  \centering
  \includegraphics[width=\columnwidth]{figures/fig_searchspace.pdf}
  \caption{Misconception and error edges in each subgraph, against the 569 in the whole KG.}
  \label{fig:searchspace}
\end{figure}

\paragraph{Larger subgraphs split into pieces.} Five of the ten subgraphs have more than one connected piece (up to four),
because the KG has no edge between some constructs, for example between output formatting and loop control in S007. A
traversal therefore needs more than one starting point.

\paragraph{The learning-progression links are in the subgraph.} Every subgraph carries its concepts' level nodes and the
links from level ways to misconceptions (69 per subgraph on average), so the rules that turn a matched misconception into
a level are available to Stage~2 without another lookup. For example, S007's subgraph contains the misconception edge
\code{loop\_control}$\to$OffByOneBound and the way that links it to \code{loop\_control@L0}.

\begin{table*}[t]
\caption{Stage~1 per scenario: concept nodes (generic kept), relation edges, misconception and error edges,
misconception nodes, level nodes, level--misconception links and disconnected pieces.}
\label{tab:perscenario}
\small
\begin{tabular}{lrrrrrrr}
\toprule
ID & Concepts (generic) & Relation edges & Mis.\ edges & Mis.\ nodes & Level nodes & Links & Pieces \\
\midrule
S000 (ref.) & 4 (2) & 2 & 18 & 17 & 20 & 20 & 2 \\
S001 & 4 (2) & 5 & 18 & 15 & 20 & 16 & 1 \\
S007 & 17 (5) & 14 & 141 & 104 & 85 & 145 & 3 \\
S008 & 19 (4) & 16 & 141 & 108 & 95 & 161 & 4 \\
S021 & 12 (3) & 10 & 87 & 69 & 60 & 91 & 2 \\
S026 & 2 (1) & 1 & 17 & 15 & 10 & 17 & 1 \\
S047 & 2 (1) & 1 & 7 & 6 & 10 & 11 & 1 \\
S052 & 8 (1) & 7 & 17 & 16 & 40 & 23 & 2 \\
S057 & 26 (5) & 22 & 167 & 120 & 130 & 193 & 4 \\
S134 & 2 (1) & 1 & 17 & 15 & 10 & 17 & 1 \\
S137 & 4 (2) & 5 & 18 & 15 & 20 & 16 & 1 \\
\bottomrule
\end{tabular}
\end{table*}

\paragraph{Limitations.} The scenarios are 10 hand-picked edge cases and one reference case, not a random sample, so the
results show feasibility, not accuracy. They were kept because their code shows the concepts the items are about; items
whose concepts appear only in the student's answer are left to Stage~2. The KG covers one chapter of one textbook, and the
reasoning stage is not built yet, so RQ2--RQ4 have no results.

\section{Discussion: What Stage~1 Implies for the Traversal}
\label{sec:discussion}
The preliminary results shape the traversal design in four ways.

\paragraph{Start from the question, not from the whole subgraph.} ToG starts from topic entities named in the question
\cite{sun2024tog}. A subgraph holds every construct in the code, up to 26 concepts. The traversal will start from the
item's main concepts and treat the rest of the subgraph as reachable context.

\paragraph{Allow several starting points.} Subgraphs with more than one piece (Table~\ref{tab:perscenario}) cannot be
covered by a walk from one entity. Following ToG, the beam will hold up to $N$ paths from different starting entities,
so that an output claim about formatting and a reason about the loop bound are explored in parallel.

\paragraph{Add the student's words as a second entry point.} A student can name a concept the code does not show,
for example a construct they compare with the one in the code. Before traversal, concepts named in the answer will be added through the aliases
stored on the concept nodes, which the extractor deliberately ignores. This keeps Stage~1 code-only and testable while
giving the traversal the concepts it needs.

\paragraph{Cost.} With the ToG bound of 22 LLM calls per clause and the two to three clauses of a typical answer, a
scenario costs at most about 66 calls, against one call for an LLM without the graph. RQ4 asks whether the gain in
accuracy and consistency justifies that cost; a RoG-style single call over the extracted paths (B2) is the cheaper
alternative we will compare against.

\section{Remaining Work and Timeline}
\label{sec:plan}
Table~\ref{tab:timeline} gives the plan to the end of the semester.

\begin{table}[t]
\caption{Week-by-week plan (Fall 2026).}
\label{tab:timeline}
\small
\begin{tabular}{p{0.16\columnwidth}p{0.74\columnwidth}}
\toprule
Week & Milestone \\
\midrule
Nov 2--8 & Clause segmentation; ToG prompts; S007 end to end. \\
Nov 9--15 & ToG traversal over Stage~1 and expert-chosen subgraphs; logging of paths and LLM calls. \\
Nov 16--22 & Learning-progression rules and root cause; RQ2 and RQ3 metrics; baseline B1 (LLM without graph). \\
Nov 23--29 & Baselines B2 (RoG-style) and B3 (lexical); ablations; five reruns for consistency (RQ4). \\
Nov 30--Dec 6 & Second annotator ($\kappa$); error analysis; answer-side concept retrieval. \\
Dec 7--13 & Final paper, slides and repository release. \\
\bottomrule
\end{tabular}
\end{table}

\paragraph{Risks and contingencies.} \emph{LLM access or cost:} the traversal is model-agnostic; if API access is
limited we run an open model locally and report the lexical judge (B3) as a floor. \emph{Small dataset:} we report
per-scenario results and bootstrap confidence intervals, and frame the study as a case study; collecting further
student answers is planned beyond this semester. \emph{Label validity:} a second annotator labels the 10 scenarios
and we report agreement. \emph{Non-determinism:} temperature~0 and five reruns per scenario. \emph{Concepts outside the
code:} retrieve concepts from the question and answer through the aliases stored in the KG before traversal.

\begin{acks}
I thank my advisor, Dr.\ Vijayalakshmi Ramasamy, for guidance. Generative AI disclosure: Claude (Anthropic) was used to
help draft and edit this report, generate figures from the measured results, and assist with code and pattern
development in the repository; all results were produced by running the repository's code, and all content was
reviewed by the author.
\end{acks}

\bibliographystyle{ACM-Reference-Format}
\bibliography{references}

\end{document}
